\documentclass[11pt]{article}
\usepackage[margin=1in]{geometry}
\usepackage{amsmath,amssymb,amsthm,mathtools}
\usepackage[colorlinks=true,linkcolor=blue,citecolor=blue,urlcolor=blue]{hyperref}
\newtheorem{theorem}{Theorem}
\newtheorem{lemma}[theorem]{Lemma}
\newtheorem{corollary}[theorem]{Corollary}
\newtheorem{proposition}[theorem]{Proposition}
\theoremstyle{definition}\newtheorem{definition}[theorem]{Definition}
\theoremstyle{remark}\newtheorem{remark}[theorem]{Remark}
\newcommand{\T}{\mathbb T_{2\pi}}
\newcommand{\C}{\mathbb C}
\newcommand{\R}{\mathbb R}
\newcommand{\norm}[1]{\left\lVert#1\right\rVert}
\newcommand{\abs}[1]{\left|#1\right|}
\title{Uniform Random Row Subsampling of Multidimensional\par Multi-Clump Vandermonde Matrices}
\author{Working mathematical draft}
\date{October 8, 2026}
\begin{document}
\maketitle
\begin{abstract}
For fixed angular nodes in dimension $d\ge1$, uniform sampling without
replacement from a cube of integer frequencies preserves the full
Vandermonde Gram matrix. The sufficient sample count is
$3072\,512^{d-1}\rho^{-2}\sum_a n_a^{2d}\log(2n/\delta)$.
The clump model is the periodic $\ell^\infty$ model used in the NumDetect
manuscript. No restriction on the internal arrangement of a clump is imposed.
The smallest singular value has a lower bound of order
$(M\Delta)^{n^\star-1}$, and this exponent is sharp over the allowed model.
An additional, generally weaker upper bound recovers the existing
one-dimensional two-sided estimate exactly when $d=1$.
\end{abstract}

\section{The model and the main theorem}
For real representatives of $x,y\in\T^d:=\R^d/(2\pi\mathbb Z)^d$, put
\[
 d_\infty(x,y)=\max_{1\le r\le d}\min_{p\in\mathbb Z}
       \abs{x_r-y_r+2\pi p},\qquad
 d_1(x,y)=\sum_{r=1}^d\min_{p\in\mathbb Z}\abs{x_r-y_r+2\pi p}.
\]
Thus $d_\infty\le d_1\le d\,d_\infty$. On representatives in
$(-\pi,\pi]^d$, each coordinate distance is
$\min\{\abs{x_r-y_r},2\pi-\abs{x_r-y_r}\}$, exactly as in
Definition \texttt{defi:metric\_separation} of NumDetect.

\begin{definition}[NumDetect clumps]\label{def:multidimensional-clumps}
Let $0<\tau\le\eta$ and $n^\star\ge2$. A set of distinct angular nodes
$\mathcal X=\{x_1,\ldots,x_n\}$ forms
$(A,\infty,\tau,\eta,n^\star)$-clumps if it admits a partition into
nonempty, pairwise disjoint sets $\mathcal C_1,\ldots,\mathcal C_A$ with
\[
 n_a:=\abs{\mathcal C_a}\le n^\star,\quad
 \max_a n_a=n^\star,\quad
 \max_{x,y\in\mathcal C_a}d_\infty(x,y)\le\tau,\quad
 \min_{x\in\mathcal C_a,\,y\in\mathcal C_b}d_\infty(x,y)>\eta\quad(a\ne b).
\]
This is Definition \texttt{defi:high\_dim\_clumps} of NumDetect. In particular,
the number of nodes at distance at most $\tau$ from a node is the size of
its clump, so its maximum is $n^\star$.
\end{definition}

Let $\mathcal K_{d,M}=\{0,\ldots,M\}^d$ and $N=(M+1)^d$.
For $\Omega\subseteq\mathcal K_{d,M}$ of cardinality $m$, define
\[
 A_M=N^{-1/2}[e^{i\langle k,x_j\rangle}]_{k\in\mathcal K_{d,M},\,1\le j\le n},
 \qquad
 A_\Omega=m^{-1/2}[e^{i\langle k,x_j\rangle}]_{k\in\Omega,\,1\le j\le n}.
\]
The order of the selected rows is immaterial. Both matrices are independent
of the representatives of their nodes. Singular values are listed in
decreasing order, with $n$ column singular values, including zeros when
necessary. Set
\[
 S_d=\sum_{a=1}^A n_a^{2d},\qquad
 C_d=3072\,512^{d-1},\qquad
 q(d,s)=\max\{q\in\mathbb N:q^d<s\}\quad(s\ge2).
\]
Here $q(d,s)\ge1$ and $q(1,s)=s-1$.

\begin{theorem}[Cube sampling and clump spectrum]
\label{thm:multidimensional-multi-clump-random-sampling}
For every $d\ge1$ and $2\le n^\star\le n$, there exist
\[
 0<c_0=c_0(d,n,n^\star)<1,\qquad C_0=C_0(d,n,n^\star)\ge n
\]
with the following property. Let $M\ge C_0$ and let the fixed distinct
angular node set $\mathcal X$ be partitioned into nonempty clumps with
$\max_a n_a=n^\star$. Suppose
\begin{equation}\label{eq:multidimensional-clump-geometry}
 \max_{x,y\in\mathcal C_a}d_\infty(x,y)\le c_0/M,\qquad
 \min_{x\in\mathcal C_a,\,y\in\mathcal C_b}d_\infty(x,y)\ge C_0/M\quad(a\ne b).
\end{equation}
Choose $\Omega$ uniformly among all $m$-element subsets of
$\mathcal K_{d,M}$. Given $0<\rho,\delta<1$ and $1\le m\le N$, assume
\begin{equation}\label{eq:multidimensional-sampling-rate}
 \boxed{m\ge\frac{C_d}{\rho^2}S_d\log\frac{2n}{\delta}.}
\end{equation}
With probability at least $1-\delta$,
\begin{equation}\label{eq:multidimensional-relative-gram}
 (1-\rho)A_M^*A_M\preceq A_\Omega^*A_\Omega
       \preceq(1+\rho)A_M^*A_M.
\end{equation}
On this same event, for every $1\le j\le n$,
\begin{equation}\label{eq:multidimensional-relative-spectrum}
 \sqrt{1-\rho}\,\sigma_j(A_M)\le\sigma_j(A_\Omega)
       \le\sqrt{1+\rho}\,\sigma_j(A_M).
\end{equation}
Moreover, for each $K\ge1$ there are positive constants $c_1,C_1$ depending
only on $d,n,n^\star,K$ such that, whenever $\Delta>0$ and
\begin{equation}\label{eq:multidimensional-comparable-spacing}
 \Delta\le d_\infty(x,y)\le K\Delta
       \quad\text{for all distinct }x,y\in\mathcal C_a,
\end{equation}
the same event satisfies
\begin{equation}\label{eq:multidimensional-sampled-smallest-singular-value}
 \boxed{c_1\sqrt{1-\rho}(M\Delta)^{n^\star-1}
       \le\sigma_{\min}(A_\Omega)
       \le C_1\sqrt{1+\rho}(M\Delta)^{q(d,n^\star)}.}
\end{equation}
The geometry thresholds are chosen before $K,\Delta,M,\mathcal X,m,\rho,\delta$.
The Gram event depends only on the fixed nodes and sampling parameters and
works simultaneously for all admissible spacing parameters and coefficient
vectors. The sample count contains no inverse separation factor.
\end{theorem}

The strict inter-clump inequality of Definition
\ref{def:multidimensional-clumps} implies the closed inequality in
\eqref{eq:multidimensional-clump-geometry} whenever
$\tau\le c_0/M$ and $\eta\ge C_0/M$. The closed boundary in the theorem
also includes the boundary of the current one-dimensional result.
There is no assumption of distinct coordinate projections, a Cartesian
product of nodes, or any particular internal arrangement.

\begin{proposition}[Sharp lower exponent]\label{prop:multidimensional-sharp-lower}
The lower estimate does not require the upper spacing inequality in
\eqref{eq:multidimensional-comparable-spacing}. Under the geometric
hypotheses of Theorem
\ref{thm:multidimensional-multi-clump-random-sampling}, a positive constant
$c=c(d,n,n^\star)$ gives
\[
 \sigma_{\min}(A_M)\ge c(M\Delta)^{n^\star-1}
\]
for every $\Delta>0$ which bounds all distinct within-clump
$d_\infty$-distances from below. Consequently, if these distances in $d_1$
are bounded below by $\Delta_1$, then
\[
 \sigma_{\min}(A_M)\ge c\,d^{-(n^\star-1)}(M\Delta_1)^{n^\star-1}.
\]
On the relative Gram event, both lower bounds acquire the factor
$\sqrt{1-\rho}$. The exponent $n^\star-1$ cannot be decreased uniformly
over the clump model, even for a single clump.
\end{proposition}

\begin{remark}[Sample count and scope]
Since $S_d\le n(n^\star)^{2d-1}$, a sufficient count for fixed $d$ and
fixed relative error is $O_{d,\rho}(n(n^\star)^{2d-1}\log(n/\delta))$.
This is a fixed-node assertion. When the sufficient count exceeds $N$,
there is no eligible cardinality; the full sample itself preserves the
Gram matrix deterministically.
\end{remark}

\section{Gap-free leverage estimates}
The following section estimates remain valid when projected nodes coincide.
This is essential in a model which imposes no internal arrangement.

\begin{lemma}[One-coordinate sections]\label{lem:angular-sections}
For fixed $n,n^\star\ge1$, there are common thresholds
$0<c_0<1$ and $C_0\ge n$ with the following properties for all
$1\le s,t\le n^\star$ and $M\ge C_0$.
If $s$ angular frequencies have lifts within $c_0/M$ of a center and
$u(k)=\sum_{j=1}^s c_j e^{ikx_j}$, then, for every
$k\in\{0,\ldots,M\}$,
\begin{equation}\label{eq:section-row-bound}
 (M+1)\abs{u(k)}^2\le512s^2\sum_{\ell=0}^M\abs{u(\ell)}^2.
\end{equation}
If two such sections have $s$ and $t$ frequencies respectively, and
their centers $x,y$ satisfy
$\abs{y-x-2\pi p}\ge C_0/M$ for every $p\in\mathbb Z$, then
\begin{equation}\label{eq:section-cross-bound}
 \abs{\langle u,v\rangle}\le\frac1{2n}\norm{u}_2\norm{v}_2.
\end{equation}
Neither assertion requires different projected frequencies.
\end{lemma}
\begin{proof}
Scale the offsets by $M$ and use the companion matrix of
$\prod_j(z-iM(x_j-\text{center}))$. Its evolution represents the exponential
sum also at repeated roots. As its roots approach zero, its first-coordinate
evolution approaches the polynomial jet basis
$1,t,\ldots,t^{s-1}/(s-1)!$ uniformly on $[0,1]$.
Polynomial evaluation and a Markov derivative estimate control the value
and the integer-grid energy. A sufficiently small perturbation gives
\eqref{eq:section-row-bound} with the displayed absolute constant.
For \eqref{eq:section-cross-bound}, approximate each section, relative to
its actual grid norm, by a modulated polynomial of degree at most $s-1$.
Finite summation by parts bounds every normalized polynomial cross moment
by $\pi/\eta$ when all winding differences of the centers are at least
$\eta$. Choose the relative approximation error first and then the
separation threshold. Finite maxima and minima supply common thresholds
for all $1\le s\le n^\star$.
\end{proof}

\begin{lemma}[Cube sections and assembly]\label{lem:cube-clump-leverage}
Thresholds depending only on $d,n,n^\star$ can be chosen so that every
single-clump signal on $\mathcal K_{d,M}$ satisfies
\[
 N\abs{u(k)}^2\le512^d n_a^{2d}\norm{u}_2^2.
\]
For different clumps their signal spaces satisfy
$\abs{\langle u,v\rangle}\le(2n)^{-1}\norm{u}_2\norm{v}_2$.
Consequently, for every full multiclump signal $f=\sum_a u_a$,
\begin{equation}\label{eq:cube-leverage}
 \sup_{f\ne0}\frac{\abs{f(k)}^2}{\norm{f}_2^2}
       \le\frac{2\,512^d S_d}{N}.
\end{equation}
\end{lemma}
\begin{proof}
Fix all coordinates except one. The resulting section is a univariate
exponential sum supported on that coordinate projection of the clump;
collisions simply repeat roots. Apply \eqref{eq:section-row-bound} in
each coordinate and sum successively. This gives the tensor factor
$(512n_a^2)^d$.
For two clumps, choose anchors. Their $d_\infty$-separation is attained in
some coordinate. Apply \eqref{eq:section-cross-bound} to the corresponding
sections and sum over all other coordinates; Cauchy--Schwarz gives the
same bound for the entire cube. The chosen coordinate can depend on the
pair of clumps.
Expanding the norm and summing the pairwise estimates gives
\[
 \norm{\sum_a u_a}_2^2\ge\tfrac12\sum_a\norm{u_a}_2^2.
\]
Weighted Cauchy--Schwarz at a fixed row gives
$N\abs{\sum_a u_a(k)}^2\le512^d S_d\sum_a\norm{u_a}_2^2$,
which proves \eqref{eq:cube-leverage}.
\end{proof}

\section{Deterministic spectrum and sharpness}
\begin{lemma}[Integer phase separation]\label{lem:integer-phase-separation}
Suppose $\eta>0$, $Q\in\mathbb N$ with $Q\ge1$, and
$\abs{\theta-2\pi p}\ge\eta$ for every integer $p$. If
$(Q+1)\eta\ge2\pi$, there exists an integer $1\le q\le Q$ with
$\abs{e^{iq\theta}-1}\ge1$.
For a short representative with $\abs{Q\theta}\le\pi$, one also has
$\abs{e^{iQ\theta}-1}\ge(2/\pi)Q\abs{\theta}$.
\end{lemma}
\begin{proof}
The geometric-sum estimate gives
$\abs{\sum_{k=0}^Q e^{ik\theta}}\le\pi/\eta\le(Q+1)/2$.
If every chord were smaller than one, the real part of each summand would
be greater than $1/2$, contradicting this bound. The second assertion is
the sine chord estimate on the principal interval.
\end{proof}

\begin{lemma}[Cardinal polynomials on a cube]\label{lem:cube-cardinal-lower}
For sufficiently large thresholds independent of the within-clump minimum
gap, the normalized full cube matrix is injective and obeys the lower
estimate in Proposition \ref{prop:multidimensional-sharp-lower}.
\end{lemma}
\begin{proof}
Fix a node $x_j$. For each other node choose a coordinate attaining its
periodic $\ell^\infty$-distance from $x_j$. Allocate an integer bandwidth
of order $M/n$ to each pair. For a pair in the same clump, the second part
of Lemma \ref{lem:integer-phase-separation} provides a chord of size at
least a constant times $M\Delta/n$. For a pair in another clump, its first
part provides a chord of size at least one.
Multiply the associated normalized two-term trigonometric factors. Scale
each within-clump factor by $\lambda=c_nM\Delta$ so that its coefficient
mass is bounded by an absolute constant. The resulting polynomial vanishes
at every other node and takes value $\lambda^{n_a-1}$ at $x_j$. Its
coefficient mass is bounded in terms of $n$, and its frequencies lie in a
cube of side at most $M/2$.
Multiply by a normalized product of Dirichlet averages in the unused
bandwidth, centered at $x_j$. Convolution bounds its coefficient norm by
its earlier coefficient mass divided by a constant times $(M+1)^{d/2}$.
Testing a Fourier signal against these $n$ coefficient vectors and using
Cauchy--Schwarz gives the claimed normalized lower bound. Since
$M\Delta\le c_0<1$ and $n_a\le n^\star$, the power $n_a-1$ may be
replaced by $n^\star-1$. A finite set of distinct within-clump distances
has a positive minimum, so this argument also proves injectivity without
any prescribed lower spacing hypothesis. The $d_1$ conclusion follows
from $d_1\le d\,d_\infty$.
\end{proof}

\begin{lemma}[Box-moment upper bound]\label{lem:box-moment-upper}
Under the geometric hypotheses of Theorem
\ref{thm:multidimensional-multi-clump-random-sampling} and
\eqref{eq:multidimensional-comparable-spacing},
$\sigma_{\min}(A_M)\le C(d,n^\star,K)(M\Delta)^{q(d,n^\star)}$.
\end{lemma}
\begin{proof}
Use a clump of size $s=n^\star$, choose short coordinate lifts relative
to an anchor, and put $q=q(d,s)$. The $q^d\times s$ matrix of box moments
indexed by $\alpha\in\{0,\ldots,q-1\}^d$ has a unit kernel vector because
$q^d<s$. This vector annihilates the tensor product of the coordinate
Taylor polynomials of degree $q-1$.
Each coordinate phase is at most $M\,\operatorname{diam}_\infty\mathcal C_a
\le c_0<1$ in absolute value. The exponential remainder and a telescoping
product estimate bound the remaining row value by
$2\sqrt{s}\,d\,3^d(MK\Delta)^q/q!$. Sum its square over the cube, divide
by $N$, and extend the vector by zero outside this clump. The variational
characterization of the least singular value gives the result.
\end{proof}

\begin{proof}[Sharpness in Proposition \ref{prop:multidimensional-sharp-lower}]
For $s=n^\star$, take the admissible single clump
$\mathcal X_\Delta=\{j\Delta e_1:0\le j\le s-1\}$ with
$(s-1)M\Delta\le c_0$. Its within-clump distances are between $\Delta$
and $(s-1)\Delta$, both for $d_1$ and $d_\infty$.
There is a unit coefficient vector annihilating the $s-1$ scalar moments
$1,j\Delta,\ldots,(j\Delta)^{s-2}$. Every Fourier phase of this collinear
family depends only on $k_1$. The univariate Taylor remainder therefore
gives, for every cube row, the upper bound
\[
 \abs{\sum_j a_j e^{i\langle k,j\Delta e_1\rangle}}
 \le B_s(M\Delta)^{s-1},\qquad
 B_s=\frac{6\sqrt{s}(s-1)^{s-1}}{(s-1)!}.
\]
After squaring and averaging, this bounds both the normalized full matrix
and every nonempty selected-row matrix by $B_s(M\Delta)^{s-1}$ in least
singular value. For any real $p<s-1$ and any $c>0$, an arbitrarily small
positive scale $t=M\Delta$ has $B_st^{s-1}<ct^p$. This contradicts a
proposed uniform lower bound of exponent $p$, while the family continues
to satisfy the literal clump model. Thus the lower exponent is sharp.
\end{proof}

\section{Sampling and exact one-dimensional specialization}
\begin{proof}[Proof of Theorem
\ref{thm:multidimensional-multi-clump-random-sampling}]
Choose common thresholds for Lemmas \ref{lem:cube-clump-leverage} and
\ref{lem:cube-cardinal-lower}. Let $G=A_M^*A_M$, which is positive definite,
and whiten the raw row vectors by $G^{-1/2}$. Their population average
outer product is the identity. Equation \eqref{eq:cube-leverage} bounds
each whitened outer product by
$R I$, where $R=2\,512^d S_d$.
The finite matrix Chernoff bound for uniform sampling without replacement
has failure probability at most $2n\exp(-m\rho^2/(3R))$.
Since $3R=3072\,512^{d-1}S_d$, condition
\eqref{eq:multidimensional-sampling-rate} makes this at most $\delta$.
Undo whitening to obtain \eqref{eq:multidimensional-relative-gram}.
Eigenvalue monotonicity for Hermitian matrices and the nonnegative square
root give \eqref{eq:multidimensional-relative-spectrum}.
The deterministic lower and upper lemmas supply the last assertion on
this same event, simultaneously for all admissible $K,\Delta$.
\end{proof}

\begin{corollary}[The current one-dimensional theorem]
\label{cor:one-dimensional-multi-clump-random-sampling}
For $d=1$, Theorem
\ref{thm:multidimensional-multi-clump-random-sampling} is exactly the current
one-dimensional random-retained-row multiclump theorem: $N=M+1$,
$S_1=\sum_a n_a^2$, the sampling constant is $3072$,
$q(1,n^\star)=n^\star-1$, and
$d_\infty=d_1=d_{\mathbb T}$. Its geometry thresholds depend only on $n,n^\star$,
and the comparable-spacing spectral constants depend only on
$n,n^\star,K$. In particular,
\[
 m\ge\frac{3072}{\rho^2}\Bigl(\sum_a n_a^2\Bigr)\log\frac{2n}{\delta},
 \qquad
 c_1\sqrt{1-\rho}(M\Delta)^{n^\star-1}
 \le\sigma_{\min}(A_\Omega)
 \le C_1\sqrt{1+\rho}(M\Delta)^{n^\star-1}.
\]
\end{corollary}
\begin{proof}
Identify a one-coordinate frequency function with its coordinate. This is
a bijection from $\{0,\ldots,M\}^1$ to $\{0,\ldots,M\}$ and induces a
bijection between their $m$-element subsets, preserving the uniform law.
Under this identification the full and sampled matrices agree up to a
row permutation. Their Gram matrices and all singular values therefore
agree. Substitute the displayed identities into the theorem.
\end{proof}

\begin{remark}[Meaning of sharpness]
Sharpness is uniform over the admissible clump model. It does not assert
that every multidimensional configuration has the same asymptotic order.
For example, the three nodes $0,\Delta e_1,\Delta e_2$ need not have the
same least-singular-value exponent as three collinear nodes. The model
allows both arrangements, and the collinear family establishes the
worst-case exponent without adding any assumption to the theorem.
\end{remark}

\begin{thebibliography}{9}
\bibitem{bdgy}
D. Batenkov, B. Diederichs, G. Goldman, and Y. Yomdin,
\emph{The spectral properties of Vandermonde matrices with clustered nodes},
2021, \href{https://arxiv.org/abs/1909.01927}{arXiv:1909.01927}.
\bibitem{emm}
T. Erd\'elyi, C. Musco, and C. Musco,
\emph{Fourier sparse leverage scores and approximate kernel learning},
2020, \href{https://arxiv.org/abs/2006.07340}{arXiv:2006.07340}.
\bibitem{bg}
D. Batenkov and G. Goldman,
\emph{Single-exponential bounds for the smallest singular value of
Vandermonde matrices in the sub-Rayleigh regime},
2021, \href{https://arxiv.org/abs/2107.09326}{arXiv:2107.09326}.
\bibitem{tropp}
J. A. Tropp, \emph{Improved analysis of the subsampled randomized Hadamard
transform}, 2011, \href{https://arxiv.org/abs/1011.1595}{arXiv:1011.1595}.
\bibitem{grossnesme}
D. Gross and V. Nesme, \emph{Note on sampling without replacing from a
finite collection of matrices}, 2010,
\href{https://arxiv.org/abs/1001.2738}{arXiv:1001.2738}.
\bibitem{kunisnagel}
S. Kunis and D. Nagel, \emph{On the smallest singular value of multivariate
Vandermonde matrices with clustered nodes}, 2020,
\href{https://arxiv.org/abs/1907.07119}{arXiv:1907.07119}.
\end{thebibliography}
\end{document}
